\chapter{Les prisonniers}
\label{chap:les-prisonniers}

Le retour au village fut plus lent que l'aller. Les prisonniers avançaient
avec les guerriers, chargés pour certains d'une partie de ce que les Orcs
avaient décidé de rapporter. Lorsqu'ils franchirent enfin la palissade, les
habitants vinrent à leur rencontre pour découvrir le résultat de l'expédition.

Le butin fut rapidement réparti. Les armes rejoignirent celles du village, les
provisions furent mises de côté et les outils confiés à ceux qui pourraient en
avoir l'usage. Certains objets changèrent plusieurs fois de mains avant de
trouver leur place dans une habitation ou un atelier. Les prisonniers, eux,
furent conduits dans un espace fermé par une simple clôture de gros rondins,
sous la surveillance de quelques guerriers.

Khaer Ghal aurait volontiers continué à accompagner les patrouilles dans les
jours qui suivirent. Maintenant qu'il avait découvert ce qui existait au-delà
des environs immédiats du village, les tâches ordinaires lui semblaient
beaucoup moins intéressantes. On lui rappela cependant qu'accompagner les
adultes impliquait également de participer aux corvées qui découlaient de
leurs expéditions.

On lui confia donc une nouvelle tâche : apporter à manger aux prisonniers.

Les consignes étaient simples. Deux fois par jour, Khaer Ghal devait récupérer
des récipients contenant une ration sommaire et suffisamment d'eau pour que les
captifs puissent travailler. Il les transportait jusqu'à la clôture, distribuait
leur contenu sous la surveillance d'un adulte, récupérait ce qui devait l'être
et repartait. Il n'était pas là pour discuter. Si quelqu'un cherchait à
l'interpeller, il devait simplement lui demander de reculer et poursuivre sa
corvée.

Les premiers jours, cette règle ne lui posa aucun problème.

Certains prisonniers tentaient bien de lui parler lorsqu'il approchait. Ils
réclamaient davantage d'eau, demandaient ce qu'on comptait faire d'eux ou
voulaient savoir combien de temps ils resteraient là. Khaer Ghal répétait ce
qu'on lui avait dit de répondre, ou ne répondait pas du tout. Il avait déjà vu
des prisonniers au village et ceux-là ne lui semblaient pas particulièrement
différents.

Puis, de temps à autre, l'un d'eux disparaissait.

Khaer Ghal savait parfaitement où ils allaient. Lorsqu'un combat était organisé
dans l'arène, un prisonnier pouvait être choisi pour affronter un guerrier du
village ou l'une des créatures capturées dans les environs. Certains
revenaient, souvent blessés. D'autres non. Là encore, il n'y voyait rien qu'il
ne connût déjà.

Parmi les nouveaux arrivants se trouvait cependant un garçon de son âge.

Khaer Ghal l'avait remarqué dès le premier jour, sans véritablement s'y
intéresser. L'enfant restait généralement près de ses parents et semblait en
meilleur état que certains adultes après le voyage jusqu'au village. Il avait
les cheveux sombres, légèrement désordonnés, et observait beaucoup ce qui se
passait autour de lui.

Surtout, il observait Khaer Ghal.

Lors de l'une des distributions, alors que celui-ci déposait un récipient près
de la clôture, le garçon s'approcha.

— C'est quoi, à ton oreille ?

Khaer Ghal releva les yeux.

L'humain regardait la longue dent qui traversait son oreille gauche, juste sous
le bord irrégulier laissé par la morsure de sa première épreuve.

— Recule.

Le garçon obéit.

Khaer Ghal termina sa distribution et repartit.

Le lendemain, il recommença.

— Ca vient d'un animal ?

— Recule.

Cette fois encore, l'enfant obéit, mais Khaer Ghal surprit son regard posé sur
la dent pendant qu'il s'éloignait.

Le jour suivant, la question changea.

— Tu as quel âge ?

Khaer Ghal posa un bol.

— Recule.

— Moi, j'en ai dix.

Khaer Ghal tourna la tête vers lui.

Il n'avait rien demandé.

Le garçon lui adressa simplement un regard curieux avant de retourner auprès
de ses parents.

Les jours suivants, il continua ainsi. Il ne protestait jamais et ne demandait
ni à sortir ni à recevoir davantage de nourriture. Il posait simplement une
question, reculait lorsqu'on le lui ordonnait et recommençait lors de la
distribution suivante.

Khaer Ghal finit par remarquer qu'il attendait presque de découvrir quelle
serait la prochaine.

Il n'aurait pas su expliquer pourquoi cet humain-là lui paraissait différent
des autres. Peut-être parce qu'il ne s'adressait pas à lui comme un prisonnier
s'adressait à l'un de ses gardiens. Il avait surtout l'air d'un enfant curieux
qui cherchait à comprendre un autre enfant.

Et Khaer Ghal connaissait assez bien la curiosité pour reconnaître ce qu'elle
pouvait provoquer.

Les parents du garçon appréciaient beaucoup moins ces échanges. Sa mère le
rappelait régulièrement lorsqu'il s'approchait trop près de la clôture et son
père lui demandait parfois de cesser de parler aux Orcs. Ils avaient assisté à
l'attaque de leur caravane et savaient ce dont les guerriers du village étaient
capables.

Mais Khaer Ghal n'avait jamais menacé leur fils.

Les autres gardiens finirent eux aussi par tolérer ces quelques paroles. Le
garçon obéissait lorsqu'on lui disait de reculer, et Khaer Ghal accomplissait
toujours sa tâche. Tant qu'aucun des deux ne causait de problème, personne ne
voyait vraiment de raison d'intervenir.

Quelques jours passèrent ainsi.

Un matin, alors que Khaer Ghal arrivait avec les rations, le garçon s'approcha
comme à son habitude. Khaer Ghal lui tendit l'un des bols à travers l'espace
entre deux rondins et l'humain avança les mains pour le recevoir.

Leurs regards se croisèrent.

%\illustration
%  {0400.png}
%  {La première rencontre}
%  {les-prisonniers}

Le garçon prit le récipient mais ne repartit pas immédiatement.

— Comment tu t'appelles ?

Khaer Ghal resta un instant silencieux. La question était différente des
précédentes. Elle ne concernait ni son oreille, ni le village, ni ce qu'il
faisait.

Elle le concernait lui.

Il aurait pu lui ordonner de reculer comme chaque fois. Au lieu de cela, il
répondit :

— Khaer Ghal.

Ce furent les premiers mots qu'il lui adressa autrement que pour lui donner un
ordre.

Le visage de l'enfant s'éclaira.

— Moi, c'est Kamatsu.

Khaer Ghal hocha légèrement la tête, puis reprit sa distribution.

Le lendemain, alors qu'il s'apprêtait déjà à repartir, une voix l'appela
derrière lui.

— Khaer Ghal !

Il s'arrêta.

Kamatsu était près de la clôture.

— Quoi ?

Le garçon sembla presque surpris d'avoir obtenu une réponse.

— Rien.

Khaer Ghal fronça les sourcils.

— Alors pourquoi tu m'appelles ?

Kamatsu réfléchit très sérieusement.

— Pour voir si tu répondrais.

Khaer Ghal le fixa quelques secondes.

— C'est idiot.

— Mais tu as répondu.

Il n'avait rien à répliquer à cela.

Khaer Ghal repartit.

\scenebreak

À partir de ce jour, leurs échanges changèrent lentement. Khaer Ghal ne devint
pas soudainement bavard et Kamatsu continua parfois à recevoir pour seule
réponse un haussement d'épaules ou un regard perplexe. Mais il ne lui ordonnait
plus systématiquement de reculer. Une question en amenait une autre, puis une
autre encore, et les quelques mots échangés pendant la distribution finirent
par devenir de véritables conversations.

Kamatsu voulait tout savoir.

Il demanda pourquoi les Orcs vivaient sur cette colline, ce qu'il y avait dans
les bois, comment Khaer Ghal avait appris à suivre une piste et pourquoi il
portait une dent à l'oreille. Cette dernière question revint suffisamment
souvent pour que Khaer Ghal finisse par lui raconter son épreuve.

— Tu l'as combattue tout seul ?

— Oui.

— À huit ans ?

— Oui.

Kamatsu regarda de nouveau son oreille.

— Elle t'a fait ça ?

Khaer Ghal porta instinctivement les doigts vers le contour abîmé du cartilage.

— Elle m'a mordu.

— Et après ?

— Je l'ai tuée.

Kamatsu resta silencieux quelques secondes.

— Et tu as gardé sa dent.

Cette fois, Khaer Ghal eut un petit mouvement de fierté.

— C'était ma première épreuve.

Kamatsu sembla réfléchir à cette réponse sans savoir exactement quoi en faire.

Les conversations reprirent les jours suivants.

Un après-midi, alors que Khaer Ghal récupérait les récipients vides, Kamatsu
désigna discrètement les autres prisonniers.

— Pourquoi vous nous gardez ici ?

Khaer Ghal haussa les épaules.

— Pour travailler.

— Et après ?

La question lui fit froncer les sourcils.

— Après quoi ?

— Quand on aura travaillé.

Khaer Ghal réfléchit. Il n'avait jamais envisagé la question de cette manière.

— Certains travailleront encore. D'autres iront dans l'arène.

Kamatsu cessa de jouer avec le bord du récipient qu'il tenait.

— Pour quoi faire ?

Khaer Ghal trouva la question étrange.

— Combattre.

— Contre quoi ?

— Ça dépend. Un guerrier. Une créature.

Kamatsu attendit la suite.

Il n'y en avait pas.

— Et si on ne veut pas ?

Khaer Ghal fronça de nouveau les sourcils.

— Pourquoi tu ne voudrais pas ?

Cette fois, ce fut Kamatsu qui le regarda comme si la réponse était évidente.

— Parce que je pourrais mourir.

Khaer Ghal désigna son oreille.

— Moi aussi, j'aurais pu mourir.

— Mais toi, tu voulais y aller ?

Khaer Ghal ouvrit la bouche, puis hésita.

— C'était mon épreuve.

— Ce n'est pas ce que j'ai demandé.

La réponse le dérangea davantage qu'elle n'aurait dû.

Khaer Ghal avait voulu réussir. Il avait voulu montrer à son père et aux autres
qu'il en était capable. Il se souvenait parfaitement de sa fierté lorsqu'il
avait enfin pu porter la dent de la créature.

Mais avait-il voulu entrer dans l'arène ?

Il chercha dans ses souvenirs le moment où quelqu'un lui aurait posé la
question.

Il n'en trouva aucun.

— C'est comme ça, finit-il par répondre.

Kamatsu baissa les yeux vers ses mains.

— Chez moi, on ne fait pas ça.

Khaer Ghal releva immédiatement la tête.

Jusque-là, Kamatsu avait surtout posé des questions sur son monde. Cette fois,
en quelques mots, il venait de rappeler qu'il en existait un autre.

Pour la première fois depuis qu'ils avaient commencé à parler, ce fut Khaer
Ghal qui posa la question.

— Alors vous faites quoi ?

Kamatsu leva les yeux vers lui.

Puis il sourit.

Et il lui raconta.

\scenebreak

À partir de ce jour, leurs conversations cessèrent progressivement de se
limiter aux distributions.

Au début, Khaer Ghal repassait simplement près de la clôture lorsqu'il avait
terminé ses tâches. Il prétendait parfois avoir oublié un récipient ou vouloir
vérifier quelque chose, même lorsque personne ne lui demandait d'explication.
Puis il cessa de chercher des excuses. Lorsque le village retrouvait son calme
en fin de journée, il venait parfois s'asseoir de son côté de la barrière et
Kamatsu s'installait en face de lui.

Le garçon avait toujours quelque chose à raconter.

Il parlait des routes qu'il avait parcourues avec ses parents, de villages où
les maisons se touchaient presque, de marchés suffisamment grands pour que
plusieurs langues s'y mêlent, de voyageurs venus de régions qu'il ne
connaissait lui-même qu'à travers leurs récits. Il décrivait des auberges
bruyantes, des ponts de pierre, des champs qui s'étendaient bien au-delà de ce
que Khaer Ghal imaginait possible et des villes entourées de murailles auprès
desquelles la palissade de son village aurait semblé minuscule.

Khaer Ghal écoutait.

Au début, il interrompait surtout pour demander des précisions.

— Combien de maisons ?

— Je ne sais pas.

— Cent ?

Kamatsu rit.

— Beaucoup plus.

— Deux cents ?

— Khaer Ghal, je ne les ai pas comptées.

Il trouva cette réponse profondément insatisfaisante.

Une autre fois, Kamatsu lui parla de la mer.

— C'est comme un lac ?

— Beaucoup plus grand.

— Jusqu'où ?

— Jusqu'à ce que tu ne voies plus rien.

Khaer Ghal le regarda avec méfiance.

— Ça n'existe pas.

— Si.

— Tu l'as vue ?

Kamatsu hésita.

— Non.

— Alors comment tu sais ?

— Mon père l'a vue.

Khaer Ghal considéra cette preuve pendant quelques secondes.

— Je demanderai à ton père.

Kamatsu éclata de rire.

Peu à peu, Khaer Ghal commença à attendre ces histoires avec autant
d'impatience qu'il avait autrefois attendu ses premières patrouilles. Certaines
lui paraissaient presque impossibles ; d'autres soulevaient immédiatement dix
nouvelles questions. Il découvrait surtout que le monde ne se limitait pas à
la colline, aux bois qui l'entouraient et aux quelques routes que les guerriers
du village surveillaient.

Il était immense.

Et plus Kamatsu lui en parlait, plus Khaer Ghal voulait savoir ce qu'il y avait
encore au-delà.

En échange, il commença à apporter les restes de son repas lorsqu'il en avait.
L'arrangement ne fut jamais réellement décidé entre eux. Kamatsu racontait une
histoire, Khaer Ghal partageait ce qu'il avait réussi à mettre de côté, puis
ils continuaient à parler jusqu'à ce que quelqu'un les rappelle.

Une histoire contre quelques morceaux de nourriture.

Cela leur convenait parfaitement.

Les parents de Kamatsu continuaient de surveiller leurs échanges avec prudence,
mais leur méfiance envers Khaer Ghal s'atténuait lentement. Ils voyaient bien
que l'Orc ne venait plus seulement accomplir une corvée et que leur fils
semblait attendre ses visites. Parfois, son père corrigeait même un détail
d'une histoire racontée un peu trop librement par Kamatsu, ce qui déclenchait
aussitôt une nouvelle série de questions de la part de Khaer Ghal.

Sans vraiment s'en rendre compte, les deux garçons apprenaient à se connaître.

Pour Kamatsu, Khaer Ghal était de moins en moins l'un des Orcs qui retenaient
sa famille au village. Pour Khaer Ghal, Kamatsu n'était déjà plus simplement
l'un des humains enfermés derrière la clôture.

Et surtout, quelque chose d'autre commençait à changer.

Les questions de Kamatsu n'apportaient pas toujours des réponses. Certaines
avaient même l'effet inverse : elles obligeaient Khaer Ghal à regarder des
choses auxquelles il n'avait jamais pensé.

Pourquoi un prisonnier devait-il travailler pour quelqu'un d'autre ?

Pourquoi devait-il combattre dans l'arène s'il ne le voulait pas ?

Pourquoi Khaer Ghal avait-il eu le droit de décider de la vie de la créature
qu'il avait affrontée, alors que personne ne lui avait jamais demandé s'il
voulait lui-même entrer dans cette fosse ?

Il ne rejetait pas encore ce qu'on lui avait appris. Il n'aurait même pas su
mettre un nom sur le malaise qui apparaissait parfois lorsque Kamatsu posait
l'une de ces questions.

Mais pour la première fois, certaines réponses qu'il avait toujours considérées
comme évidentes ne l'étaient plus tout à fait.

Une dizaine de jours s'était écoulée depuis l'arrivée des prisonniers,
peut-être davantage. Khaer Ghal avait cessé de compter.

Il savait seulement qu'il reviendrait le lendemain.

Et qu'il aurait probablement encore des questions.